\section{Discussion}
The results support a focused methodological conclusion: prompt objectives can be treated as controlled interventions, and SBERT orientation provides an interpretable measure for ordering, selecting, and evaluating their semantic direction. Holding the market environment, action space, prompt structure, and model configuration fixed, more profit-oriented objective blocks generally produce higher markups, prices, and firm profits.

Competitive market and Nash band benchmarks anchor this semantic result in economic terms. Figure~\ref{fig:price_trajectories} shows that price elevation is concentrated in peak and near-scarcity hours, where strategic markups have the greatest opportunity to affect clearing prices. The perfect competition benchmark provides the cost-based reference path, while the Nash band provides an equilibrium reference point under the finite markup grid. As prompts become more profit-oriented, agent bids move market outcomes away from the perfect competition benchmark and closer to the Nash benchmark band. This does not mean that the agents solved for or converged to Nash equilibrium. Rather, it shows that profit-oriented objectives induce more aggressive bidding behavior that moves outcomes closer to the highest-profit strategic benchmark available in the simplified action space.

The SBERT orientation regression coefficients provide quantitative evidence for this trend. A 0.1 increase in prompt orientation is associated with higher price impacts, markups, firm profit gains, and consumer payment increases. These estimates are useful because they quantify the relationship between prompt orientation and economic units. These estimates should be interpreted as descriptive evidence from a controlled case study rather than as large-sample statistical inference. The relationship is estimated from six treatment-level observations, so the result is informative about the prompt interventions studied here but not yet sufficient to establish a general law across all possible prompt treatments.

The semantic results also show that prompt orientation and generated reasoning orientation move in the same direction. Both are aligned with the agents' bidding posture. We find that more profit-oriented objectives produce higher markups and prices, and more profit-oriented TARJ language accompanies more aggressive bidding behavior. This supports the use of semantic orientation not only as a prompt-design measure, but also as a diagnostic for whether the agent's generated reasoning artifacts are directionally consistent with its actions.

